% Proposed method: Pf_Filter
% Needs in the preamble: \usepackage{graphicx}, \usepackage{booktabs} (optional), \usepackage{amsmath}

\section{Proposed Method: Pf\_Filter}
\label{sec:pf-filter}

Pf\_Filter is a software prefetching scheme for the HNSW base-layer search.
It requests each piece of data a few steps before the search needs it, and it
uses the cheap visited flag to skip fetching vectors that the search will discard.
It changes only \emph{when} data is loaded, never \emph{what} the search computes.

\subsection{Where a Search Step Waits}

Each step of the base-layer search does the following:
\begin{enumerate}
  \item Pop the closest candidate from the priority queue.
  \item Read its neighbour list.
  \item For each neighbour: read its visited flag; if it is new, load its vector,
        compute the distance and, if it is close enough, push it into the queue.
\end{enumerate}
The neighbour list, the visited flag and the vector each sit at an effectively
random address. Each one is a separate trip to DRAM costing a few hundred cycles.
This is the idle time measured in Section~\ref{sec:observation}.

\subsection{Why the Stock Prefetch Falls Short}

Stock hnswlib prefetches the visited flag and the first cache line of neighbour
$j+1$ while it works on neighbour $j$. This has two weaknesses:
\begin{itemize}
  \item \textbf{Too late.} Processing one neighbour takes less time than one DRAM
        access, so the data for $j+1$ has not arrived when the loop reaches it.
  \item \textbf{Too little.} One 64-byte line covers a small part of a vector:
        a SIFT vector is 512\,B, a GIST vector 3{,}840\,B and an Enron vector 5{,}476\,B.
\end{itemize}

\subsection{Pf\_Filter Design}

Pf\_Filter has three parts, with a look-ahead distance $D = 2$
(Figure~\ref{fig:pf-filter-design}).

\begin{figure}[t]
  \centering
  \includegraphics[width=\linewidth]{new_plot/pf_filter_design}
  \caption{Pf\_Filter design: one step of the HNSW base-layer loop at neighbour $j$
           ($D = 2$). At $j$ the vector is already in cache; the flag of $j+2$
           decides whether its vector is fetched, and the flag of $j+4$ is requested next.}
  \label{fig:pf-filter-design}
\end{figure}

\paragraph{1. Two-stage, filtered look-ahead.}
While processing neighbour $j$, the loop issues two prefetches:
\begin{itemize}
  \item the visited flag of neighbour $j + 2D$ ($j+4$). A flag is one byte, so
        fetching it far ahead is cheap;
  \item the first three cache lines of the vector of neighbour $j + D$ ($j+2$),
        but \emph{only if its flag says it is not yet visited}. That flag was
        requested two steps earlier and is already in cache, so the check costs
        almost nothing.
\end{itemize}
The filter is the key change. Many neighbours are already visited and will be
skipped. An unfiltered look-ahead fetches their vectors anyway, wasting memory
bandwidth and evicting useful lines. Pf\_Filter fetches only the vectors the
search will actually use.

\paragraph{2. Warm-up.}
At the start of each neighbour list nothing has been prefetched yet. The loop
first requests the flags of neighbours $1$ to $2D$ and the vectors of neighbours
$1$ to $D$, so the look-ahead is full from the first neighbour.

\paragraph{3. Next-candidate neighbour list.}
While the current node's neighbours are processed, Pf\_Filter prefetches the
neighbour list of the next candidate in the queue. A node's neighbour list is
also prefetched when it is pushed into the queue. The next step's neighbour-list
read then hits in cache.

Pf\_Filter also includes a shortcut-jump prefetch that applies only to SHG's
shortcut edges. HNSW has no shortcuts, so this part never runs in any HNSW
result in this report.

\subsection{Correctness and Cost}

A prefetch is a hint to the hardware: it moves data closer to the core but does
not change program state. The search visits the same nodes, computes the same
distances and returns the same results. We verified this: 0 mismatches against
vanilla HNSW on all four datasets and all efSearch values.

The added work per neighbour is one flag read from cache and at most two
prefetch instructions. No change is made to the index, its memory layout or
its build.
